\documentclass[12pt]{article}
\usepackage[T1]{fontenc}
\usepackage[utf8]{inputenc}
\usepackage[brazil]{babel}
\usepackage[margin=1in]{geometry}
\setlength{\parindent}{0pt}
\begin{document}
\section*{Tema 18}
Processo: PEDILEF 2009.71.95.000509-1/ RS\\
Situacao: Julgado (Súmulas 5 e 14 da TNU)\\
Ramo: DIREITO PREVIDENCIÁRIO\\
Questao: Saber se documentos em nome de terceiros integrantes do grupo familiar, relativos à propriedade da terra trabalhada, servem como início de prova material da atividade rural.\\
Tese: A certidão do INCRA ou outro documento que comprove propriedade de imóvel em nome de integrantes do grupo familiar do segurado é razoável início de prova material da condição de segurado especial para fins de aposentadoria rural por idade, inclusive dos períodos trabalhados a partir dos 12 anos de idade, antes da publicação da Lei n. 8.213/91. Desnecessidade de comprovação de todo o período de carência.\\
Fonte: https://www.cjf.jus.br/phpdoc/virtus/
\end{document}
